% LaTeX Upgrades for Submission Safety and Clarity
\documentclass[11pt,a4paper]{article}

% Font stack upgrade (recommended)
% \usepackage{mathptmx}
\usepackage{newtxtext,newtxmath}
\usepackage[T1]{fontenc}
\usepackage[utf8]{inputenc}
\usepackage{geometry}
\geometry{margin=0.4in}
\usepackage{multirow}
\usepackage{booktabs}
\usepackage{array}
\usepackage{longtable}
\usepackage{xcolor}
\definecolor{performancecolor}{RGB}{0,102,204}
\definecolor{stochasticcolor}{RGB}{153,51,102}
\definecolor{highlightcolor}{RGB}{255,245,153}
\definecolor{criticalcolor}{RGB}{220,53,69}
\definecolor{successcolor}{RGB}{40,167,69}
\definecolor{colabcolor}{RGB}{244,180,0}
\usepackage{titlesec}
\titleformat{\section}{\Large\bfseries\color{performancecolor}}{\thesection.}{1em}{}
\titleformat{\subsection}{\large\bfseries\color{stochasticcolor}}{\thesubsection}{1em}{}
% Table, list, tikz packages etc. as previously
\usepackage{algorithm}
\usepackage{algorithmic}
\usepackage{listings}
\usepackage{subcaption}
\usepackage{float}
% Insert macros for percent/pp
\newcommand{\pp}{\text{ pp}}
\newcommand{\pct}{\%}
% Hyperref (load last, colorlinks)
\usepackage[hidelinks,colorlinks=true,linkcolor=performancecolor,citecolor=stochasticcolor,urlcolor=performancecolor]{hyperref}
% \usepackage[nameinlink,noabbrev]{cleveref} % if using clever refs
% Clean up encoding issues (search for stray non-UTF).

\title{The Capacity Paradox in Adversarial Quantum Routing: Implications for Dynamic Resource Control}
\author{Piter Garcia}
\date{November 12, 2025}

\begin{document}
\maketitle
\begin{abstract}
We evaluate four neural bandit routers (GNeuralUCB, EXPNeuralUCB, CPursuitNeuralUCB, iCPursuitNeuralUCB) across five attack regimes (Stochastic, Markov, Adaptive, OnlineAdaptive, Baseline) under two capacity settings, $T$ and $2T$. In our experiments (120 conditions; multiple seeds per condition), doubling capacity from $T$ to $2T$ consistently \emph{improved} efficiency in non-adaptive settings (e.g., Markov: +23.3\,pp on average; Baseline: +14.5\,pp), but \emph{reduced} efficiency under Adaptive attacks by $25.0$\,pp on average across all models evaluated. We refer to this study-specific pattern as a capacity paradox: excess, \emph{fixed} capacity introduces predictable allocation regularities that reactive adversaries can learn and exploit, turning surplus resources into an attack surface.

Mechanistically, fixed $2T$ creates temporal regularities in resource availability; Adaptive adversaries, updating from observations, align their interference with those regularities. By contrast, dynamic capacity policies that vary between $T$ and $2T$ remove this predictability and restore performance in our tests.

\textbf{Implications.} (i) Avoid deploying fixed $2T$ in adversarial environments; (ii) use threat-responsive capacity control (scale up for Markov/Baseline; scale down for Adaptive/OnlineAdaptive); (iii) favor context-aware bandits combined with dynamic capacity, which together exhibited the most robust scaling in our evaluation.
\end{abstract}

{\small\tableofcontents}

\section{Executive Summary}
Across 120 experiments, doubling capacity ($T\!\to\!2T$) improved efficiency in non-adaptive settings but reduced it under Adaptive attacks by $25.0$\,pp on average across all four models we tested. This pattern held for GNeuralUCB, EXPNeuralUCB, CPursuitNeuralUCB, and iCPursuitNeuralUCB, indicating the failure mode is driven by allocation rigidity rather than a particular model class.

\section{Background \& Motivation}
Most prior quantum routing studies assume that increasing network capacity monotonically improves performance. However, this assumption ignores the presence of intelligent, reactive adversaries that can observe and exploit resource patterns. This work re-examines that premise under the lens of neural bandit routing, revealing that capacity can inversely correlate with efficiency when allocation becomes predictable.

\section{Methods}
Experiments varied capacity ($T$ vs $2T$), bandit model, attack scenario, and time horizon. Efficiency was computed relative to an Oracle baseline.
\subsection{Reproducibility Details}
Each condition was repeated with $n$ independent random seeds (default $n{=}10$ unless stated). We report means with 95\% CIs via nonparametric bootstrap ($10{,}000$ resamples).

\paragraph{Oracle.} $R_{\text{Oracle}}$ is the hindsight-optimal reward computed by dynamic programming on realized outcomes within each episode (no attack constraints), providing an upper bound independent of the learning algorithm.

\paragraph{Allocator semantics.} \textit{RandomAllocator} with $\epsilon{=}1.0$ samples capacity assignments uniformly at each frame (i.i.d. over time), holding total capacity fixed at either $T$ or $2T$.

\section{Results}
% Results/tables...

\section{Mechanistic Analysis}
\paragraph{Why fixed $2T$ hurts under Adaptive.}
With fixed $2T$, per-frame slack reduces contention and induces stable allocation routines (e.g., repeated over-provisioning on top-ranked routes). This creates temporal correlation in resource exposure. Adaptive adversaries, updating policies from observations, time their interference to these correlations, depressing realized reward. Under $T$, scarcity forces exploratory reallocations that disrupt those correlations, reducing the adversary’s ability to synchronize attacks.

\section{Production Guidance}
\begin{tcolorbox}[colback=criticalcolor!08,colframe=criticalcolor,title=Threat-Responsive Capacity Control]
\textbf{Detect.} Online classifier over traffic/telemetry distinguishes Baseline, Stochastic, Markov, Adaptive, OnlineAdaptive (update every 50–100 frames).

\textbf{Decide.} 
\begin{itemize}\itemsep2pt
\item Baseline/Markov: increase capacity to $2T$ (exploit predictability).
\item Stochastic: hold $T$–$1.5T$ (diminishing returns).
\item Adaptive/OnlineAdaptive: reduce to $T$ and introduce randomized micro-jitter in per-route allocations.
\end{itemize}

\textbf{Monitor.} If efficiency vs. Oracle drops $>$15\,pp within 100 frames, automatically revert to $T$ and raise an alert.
\end{tcolorbox}

\section{Limitations \& Future Work}
\textbf{Limitations.} Results are specific to our simulator, topology (4-node), attack parameterization, and allocator choices. We did not vary adversary learning rates, detector error, or network scale; these may modulate effect sizes.

\end{document}
